\documentclass[a4paper, serif, 10pt]{moderncv}

%% moderncv
% color and style
\moderncvcolor{black}
\moderncvstyle{banking}

%% page layout/margins
\usepackage[top=2.5cm, bottom=2.5cm, left=1.5cm, right=1.5cm]{geometry}

\nopagenumbers{}

%% Babel config for Spanish language
% \usepackage[spanish]{babel} has some conflicting issues with moderncv
% so we have to use enable the following options to make it work
%
% ref:
% - https://tex.stackexchange.com/a/140161/36007
\usepackage[spanish,es-lcroman]{babel}

%% fontspec
\usepackage{fontspec}

\IfFontExistsTF{Linux Libertine}{
  \setmainfont[Ligatures={TeX, Common}, Numbers=OldStyle]{Linux Libertine}
}{}
\IfFontExistsTF{Linux Libertine O}{
  \setmainfont[Ligatures={TeX, Common}, Numbers=OldStyle]{Linux Libertine O}
}{}

%% CTeX
% CJK support, used to show CJK characters in the resume
%
% - heading=false: disable ctex heading
% - punct=kaiming: use kaiming punctuations styles for CJK
% - scheme=plain: use plain scheme, do not override \normalsize font size
% - space=auto: space settings for CJK characters
%
% CJK fonts are resolved per locale language, e.g. Noto Serif/Sans CJK JP
% for Japanese resumes, with fallback to the next region if not installed
%
% ref:
% - http://ctan.mirrorcatalogs.com/language/chinese/ctex/ctex.pdf
\usepackage[UTF8, heading=false, punct=kaiming, scheme=plain, space=auto]{ctex}

\IfFontExistsTF{Noto Serif CJK SC}{
  \setCJKmainfont{Noto Serif CJK SC}
}{}
\IfFontExistsTF{Noto Sans CJK SC}{
  \setCJKsansfont{Noto Sans CJK SC}
}{}

%% line spacing
\usepackage{setspace}
\setstretch{1.125}

%% URL styling - use normal text instead of monospace
\urlstyle{same}

%% Patch \httplink and \httpslink to handle full URLs
% The original moderncv macros always prepend http:// or https:// to URLs.
% This causes issues when the URL already has a protocol (e.g., https://example.com)
% resulting in malformed URLs like https://https://example.com.
% This patch checks if the URL already contains "://" and uses it directly if so.
% Note: We use \str_set:Nx (x-expansion) to expand macros like \@homepage first.
\ExplSyntaxOn
\str_new:N \g_yamlresume_url_str

\RenewDocumentCommand{\httpslink}{O{}m}{
  \str_set:Nx \g_yamlresume_url_str {#2}
  \str_if_in:NnTF \g_yamlresume_url_str {://}
    { \url{#2} }
    { \url{https://#2} }
}

\RenewDocumentCommand{\httplink}{O{}m}{
  \str_set:Nx \g_yamlresume_url_str {#2}
  \str_if_in:NnTF \g_yamlresume_url_str {://}
    { \url{#2} }
    { \url{http://#2} }
}
\ExplSyntaxOff

\name{Lucía Fernández Moreno}{}
\title{Ingeniera de Datos Senior}
\phone[mobile]{+34 612 345 678}
\email{lucia.fernandez@example.com}
\homepage{https://luciafernandez.dev}

\address{Calle de Alcalá 123, Madrid, Comunidad de Madrid, España, 28009}{}{}

\extrainfo{{\small \faLinkedin}\ \href{https://www.linkedin.com/in/lucia-fernandez-data}{@lucia-fernandez-data} {} {} {} • {} {} {} 
{\small \faGithub}\ \href{https://github.com/luciafdez}{@luciafdez}}

\begin{document}

\maketitle

\section{Información básica}

\cvline{}{\begin{itemize}
\item Ingeniera de Datos con más de 10 años de experiencia diseñando y construyendo plataformas de datos escalables.
\item Especializada en arquitecturas de datos modernas, pipelines ETL/ELT y soluciones analíticas en la nube.
\item Sólida experiencia en modelado dimensional, optimización de consultas y gobierno del dato.
\item Colaboradora habitual con equipos de negocio, analítica y ciencia de datos para convertir requisitos en soluciones medibles.
\end{itemize}}
\section{Educación}

\cventry{sept 2012–jul 2014}
        {Maestría, Ciencia de Datos e Ingeniería de Datos, Puntuación: 9.1}
        {Universidad Politécnica de Madrid}
        {\url{https://www.upm.es}}
        {}
        {\begin{itemize}
\item Formación avanzada en diseño de sistemas de datos, analítica y aprendizaje automático.
\item Proyecto fin de máster centrado en la optimización de pipelines de ingestión en tiempo real.
\end{itemize}
\textbf{Cursos}: Arquitecturas de Datos Escalables, Procesamiento Distribuido, Minería de Datos, Bases de Datos Avanzadas, Aprendizaje Automático}

\cventry{sept 2008–jun 2012}
        {Licenciatura, Ingeniería Informática, Puntuación: 8.7}
        {Universidad de Sevilla}
        {\url{https://www.us.es}}
        {}
        {\begin{itemize}
\item Base sólida en programación, sistemas distribuidos y fundamentos de bases de datos.
\end{itemize}
\textbf{Cursos}: Estructuras de Datos, Algoritmos, Sistemas Operativos, Redes de Computadores, Bases de Datos}
\section{Experiencia laboral}

\cventry{mar 2021 hasta la fecha}
        {Ingeniera de Datos Senior}
        {Iberia Data Solutions}
        {\url{https://www.iberiadata.example.com}}
        {}
        {\begin{itemize}
\item Lideré el diseño de una plataforma de datos en Google Cloud Platform para centralizar información de 12 áreas de negocio.
\item Implementé pipelines ELT con dbt, BigQuery y Airflow, reduciendo el tiempo de publicación de datos en un 60 \%.
\item Definí estándares de modelado dimensional y gobierno del dato, mejorando la calidad y trazabilidad de los datasets.
\item Mentoré a tres ingenieros junior y coordiné revisiones técnicas de código y arquitectura.
\end{itemize}
\textbf{Palabras clave}: Google Cloud Platform, BigQuery, dbt, Airflow, Gobierno del dato}

\cventry{ene 2017–feb 2021}
        {Ingeniera de Datos}
        {Telefónica Tech}
        {\url{https://www.telefonicatech.com}}
        {}
        {\begin{itemize}
\item Desarrollé pipelines ETL en Python y Spark para procesar más de 2 TB diarios de datos de red y negocio.
\item Modelé datos en Amazon Redshift y optimicé consultas analíticas para equipos de BI y marketing.
\item Automaticé validaciones de calidad con Great Expectations y monitoreo de jobs con Grafana y Prometheus.
\item Participé en la migración de un almacén de datos on-premise a AWS con enfoque en mínima interrupción.
\end{itemize}
\textbf{Palabras clave}: Python, Apache Spark, AWS, Redshift, ETL}

\cventry{sept 2014–dic 2016}
        {Analista de Datos}
        {Accenture España}
        {\url{https://www.accenture.com/es-es}}
        {}
        {\begin{itemize}
\item Construí informes y cuadros de mando en Power BI para clientes de banca y seguros.
\item Escribí consultas SQL complejas y procedimientos almacenados para extraer y transformar datos.
\item Documenté requisitos funcionales y colaboré con equipos de negocio en la definición de KPIs.
\end{itemize}
\textbf{Palabras clave}: SQL, Power BI, Análisis de Datos, Banca}
\section{Idiomas}

\cvline{Español}{Competencia nativa o bilingüe \hfill \textbf{Palabras clave}: Español nativo}
\cvline{Inglés}{Competencia profesional plena \hfill \textbf{Palabras clave}: C1, Uso profesional}
\cvline{Catalán}{Competencia limitada de trabajo \hfill \textbf{Palabras clave}: Conversación básica}
\section{Competencias}

\cvline{Ingeniería de Datos}{Experto \hfill \textbf{Palabras clave}: ETL/ELT, Modelado dimensional, Almacén de datos, Calidad del dato, Gobierno del dato}
\cvline{Computación en la Nube}{Avanzado \hfill \textbf{Palabras clave}: Google Cloud Platform, AWS, BigQuery, Redshift, Azure Data Factory}
\cvline{Big Data y Procesamiento}{Avanzado \hfill \textbf{Palabras clave}: Apache Spark, Apache Airflow, Apache Kafka, dbt, Databricks}
\cvline{Lenguajes y Bases de Datos}{Avanzado \hfill \textbf{Palabras clave}: Python, SQL, Scala, PostgreSQL, MongoDB}
\section{Premios}

\cventry{nov 2023}
        {Iberia Data Solutions}
        {Premio a la Innovación en Datos}
        {}
        {}
        {\begin{itemize}
\item Reconocimiento por el diseño de una plataforma de datos que redujo el time-to-insight en un 45 \%.
\end{itemize}}

\cventry{jul 2014}
        {Universidad Politécnica de Madrid}
        {Matrícula de Honor en Proyecto Fin de Máster}
        {}
        {}
        {\begin{itemize}
\item Galardón al mejor proyecto de la promoción en el área de procesamiento distribuido.
\end{itemize}}
\section{Certificados}

\cventry{mar 2022}
        {Google Cloud}
        {Google Cloud Professional Data Engineer}
        {\url{https://cloud.google.com/certification/data-engineer}}
        {}
        {}

\cventry{sept 2021}
        {Amazon Web Services}
        {AWS Certified Data Analytics - Specialty}
        {\url{https://aws.amazon.com/certification/certified-data-analytics-specialty/}}
        {}
        {}

\cventry{may 2023}
        {dbt Labs}
        {dbt Analytics Engineering Certification}
        {\url{https://www.getdbt.com/certifications/}}
        {}
        {}
\section{Publicaciones}

\cventry{oct 2022}
        {Diseño de pipelines de datos en tiempo real para analítica operativa}
        {Revista Española de Ingeniería de Datos}
        {\url{https://www.revista-ingenieria-datos.example.com/articulos/pipelines-tiempo-real}}
        {}
        {\begin{itemize}
\item Artículo técnico sobre patrones de ingestión y procesamiento continuo con Kafka y Spark Structured Streaming.
\item Analiza estrategias de particionado, exactly-once y monitorización de latencia.
\end{itemize}}
\section{Proyectos}

\cventry{ene 2023–jun 2023}
        {Plataforma reproducible para ingerir, validar y publicar datasets abiertos de administraciones públicas.}
        {Observatorio de Datos Abiertos}
        {\url{https://github.com/luciafdez/observatorio-datos-abiertos}}
        {}
        {\begin{itemize}
\item Diseñé un pipeline con Airflow y dbt que ingesta datos de portales open data y genera tablas analíticas listas para consulta.
\item Implementé tests de calidad y documentación automática del linaje de datos.
\item Desplegué la solución en Google Cloud Run y BigQuery con costes controlados.
\end{itemize}
\textbf{Palabras clave}: Open Data, Airflow, dbt, BigQuery}

\cventry{mar 2020–nov 2020}
        {Prototipo de motor de recomendación basado en datos transaccionales anonimizados.}
        {Motor de Recomendación para Retail}
        {\url{https://github.com/luciafdez/retail-recommender}}
        {}
        {\begin{itemize}
\item Construí un pipeline de features con Spark y un modelo de filtrado colaborativo en Python.
\item Evalué el rendimiento con métricas de precisión y recall sobre un conjunto de validación temporal.
\item Documenté la arquitectura y los resultados en un informe técnico reproducible.
\end{itemize}
\textbf{Palabras clave}: Spark, Python, Recomendación, Retail}
\section{Intereses}

\cvline{Datos y Tecnología}{Open Source, Arquitecturas de Datos, Automatización}
\cvline{Actividad Física}{Senderismo, Natación, Ciclismo}
\cvline{Cultura}{Lectura, Fotografía, Podcasts}
\section{Voluntariado}

\cventry{sept 2022 hasta la fecha}
        {Mentora Técnica}
        {Data Girls Madrid}
        {\url{https://www.datagirlsmadrid.example.com}}
        {}
        {\begin{itemize}
\item Acompaño a mujeres que inician su carrera en datos en la definición de objetivos y desarrollo de habilidades técnicas.
\item Participo en talleres de SQL, Python y visualización de datos para la comunidad.
\end{itemize}}
    

\cventry{ene 2020–ago 2022}
        {Voluntaria de Análisis de Datos}
        {Cáritas Madrid}
        {\url{https://www.caritasmadrid.org}}
        {}
        {\begin{itemize}
\item Apoyé al equipo social en la limpieza y análisis de datos de programas de asistencia.
\item Automaticé informes mensuales en Power BI para facilitar la toma de decisiones.
\end{itemize}}
    
\end{document}